\section{Basepoint rigidity for the actual cubic deformation ring}
\label{sec:cubic-basepoint-rigidity}

Keep the actual cubic ring and graded model of AC-G49--G51,
extended by AC-G54 at the remaining split coefficient primes.
Thus $e\geq1$, $\ell\ne2,3,269$ and all retained
\texttt{GD} hypotheses hold.  Use the transported local
ideal in AC-G49-inert or the actual fixed-type prime of
AC-G54, respectively.  Write $\Gamma=G_{\mathbb Q}$,
$H=G_K$, $K=\mathbb Q(\omega)$, $\eta=\eta_K$,
$\gamma=\psi\kappa_\ell^e$ and $j\in\{1,2\}$.
The cubic seed has matrix
\[
 r_\chi(g).\mathrm{mat}
   =\operatorname{diag}\bigl(q(g)^j\gamma(g),1,
                             q(g)^j\gamma(g)^{-1}\bigr),
 \qquad r_\chi(g).\mathrm{mult}=\eta(g).
                                                               \tag{G53.1}
\]
In particular $r_\chi(c).\mathrm{mat}=I$ at the chosen
complex conjugation.  The inputs retained below are the local
published Bloch--Kato statements in \texttt{CohInp.loc},
the known strict negative-Tate finiteness used in AC-G22
and AC-K10, and ordinary class field theory.  None is inferred
from the supplied analytic theorem A.  No Kato input for the
new mixed rank-four coefficient is used here.

For a finite extension $L/\mathbb Q_\ell$ and a continuous
finite-dimensional $L$-representation $V$ of $G_{\mathbb Q}$,
write $H^1_f(\mathbb Q,V)$ for the group with unramified
conditions away from $\ell$ and the $B_{\rm cris}$ kernel
condition at $\ell$.  These are the rational local conditions
encoded by the source's \texttt{Af.IsBK}; they are not an
integral Selmer-length assertion.

\begin{lemma}[AC-G53-Artin: nonpositive Artin--Tate Bloch--Kato vanishing]
\label{AC-G53-Artin}
Let $A$ be a finite-image representation of $G_{\mathbb Q}$
over $L$.  For every integer $n\geq0$,
\[
                    H^1_f(\mathbb Q,A(-n))=0.
                                                               \tag{G53.2}
\]
The assertion permits a nonabelian splitting field and permits
its degree to be divisible by $\ell$.  It concerns rational
cohomology only.
\end{lemma}
\begin{proof}
Choose a finite Galois extension $F/\mathbb Q$ which contains
$K$ and kills the finite action on $A$.  Then $F$ is totally
imaginary, and the restriction of $A(-n)$ is a direct sum of
copies of $L(-n)$.  Restriction, coordinate projection and
finite coefficient extension preserve both local conditions:
away from $\ell$ this is immediate on inertia, and at
$\ell$ it follows by restricting a $B_{\rm cris}$ cobounding
vector and applying the coefficient projection.  Thus it
suffices to prove the corresponding vanishing over $F$ and
then use injectivity of rational restriction.

\proofstep{The exact local facts.}
For every finite place $v\mid\ell$ of $F$ and $n\geq1$,
\[
 H^1_f(F_v,L(-n))=0,
 \qquad H^1_f(F_v,L)=H^1_{\rm unr}(F_v,L).
                                                               \tag{G53.3}
\]
The first is exactly the negative-weight clause of
\path{PadicHodgeTheory/ExistsNeOne.lean},
\texttt{Loc.z1fChar}, applied to the open embedding
$G_{F_v}\hookrightarrow G_{\mathbb Q_\ell}$ and the
character $\kappa_\ell^{-n}$ of weight $-n<0$.
The second is the trivial-representation instance of
\path{PadicHodgeTheory/ExistsIntBasis.lean},
\texttt{Loc.z1fUnr}: a finite cocycle for an unramified
representation vanishes on inertia.
These statements use the retained local published inputs;
their finite-dimensional and continuity hypotheses hold
over the finite field $L$.  There is no assertion here that
a general ramified weight-zero representation is itself
unramified before passing to $F$.

\proofstep{Finite torsion strict cohomology forces rational strict vanishing.}
Suppose $n\geq1$ and put $T=\mathcal O_L(-n)$,
$V=L(-n)$ and $D_L=V/T$.
The group
\[
 \mathcal R(F,D_L)=
 \ker\!\left(H^1(G_{F,S_\ell},D_L)
       \longrightarrow\prod_{v\mid\ell}H^1(F_v,D_L)\right)
                                                               \tag{G53.4}
\]
is finite.  Indeed for $\mathbb Q_\ell/\mathbb Z_\ell(-n)$
this is the strict finiteness in AC-K10; equivalently it is
the negative-index application $i=-n$, $1-i=n+1$ of
Soul\'e, \emph{\'El\'ements cyclotomiques en K-th\'eorie},
Ast\'erisque 147--148 (1987), \S\S6.1--6.2, printed
pp.~246--247, already checked in AC-G22.
Choose a $\mathbb Z_\ell$-basis of $\mathcal O_L$.
The cyclotomic scalar action lies in $\mathbb Z_\ell$, so
this basis identifies (G53.4), as an abelian group, with
a finite direct sum of those finite groups.  This justifies
the coefficient change without an integral character
projector.

Here is a direct proof of the rational consequence that
keeps finiteness distinct from a size bound.
Let $z:G_F\to V$ be a continuous strict cocycle, using
the sign $\partial y(g)=y-\kappa_\ell(g)^{-n}y$.
Choose $\delta\in G_F$ with
$u=1-\kappa_\ell(\delta)^{-n}\ne0$; the cyclotomic image
is infinite and $n>0$.
Subtract $\partial(z(\delta)/u)$, so $z(\delta)=0$.
Compactness makes the image bounded.  After multiplying
by a nonzero scalar, assume $z$ takes values in $T$.
Both operations preserve the local strict conditions.

Let $\varpi$ be a uniformizer of $\mathcal O_L$ and choose
$a$ such that $\varpi^a\mathcal R(F,D_L)=0$.
For each $N>a$, the cocycle $\varpi^{-N}z\bmod T$
defines a member of (G53.4).  At places away from $\ell$,
$V$ is unramified, so an unramified cocycle has value zero
on inertia and factors through $G_{F,S_\ell}$.
At places above $\ell$, its class is zero by strictness.
Consequently there is $y_N\in V$ with
\[
 \varpi^{-(N-a)}z(g)
       -(1-\kappa_\ell(g)^{-n})y_N\in T
                    \qquad(g\in G_F).
                                                               \tag{G53.5}
\]
At $g=\delta$ this gives $u y_N\in T$, hence
$y_N\in u^{-1}T$ independently of $N$.
All cyclotomic values are integral units and $u\in\mathcal O_L$.
It follows from (G53.5) that
\[
                  z(g)\in\varpi^{N-a}u^{-1}T
                      \qquad\hbox{for every }N>a.
\]
The intersection of these lattices is zero.  Therefore
the normalized $z$ is zero, and its original class was
a coboundary.  By (G53.3), every Bloch--Kato class in
$H^1(F,L(-n))$ is such a strict class.  This proves the
negative-Tate case without estimating the exponent $a$.

\proofstep{Weight zero and descent.}
For $n=0$, a cocycle with trivial coefficient $L$ is
a continuous additive homomorphism.  Its local condition
away from $\ell$ and (G53.3) at $\ell$ make it zero on
all inertia groups.  Since $F$ has no real places, it
factors through the everywhere-unramified abelian
Galois group, identified by class field theory with
the finite ideal class group of $F$.  A finite group
has no nonzero homomorphism to the characteristic-zero
additive group of $L$.  This proves the zero-weight
case.  The exact same class-field step is
\texttt{unr\_zero} in
\path{UnitaryThreeGroups/BlochKatoNonpositiveWeight.lean};
its source uses \texttt{Kum.artKer} and finiteness of the
ideal class group.

Finally restriction on rational $H^1$ to a finite
extension is injective: corestriction composed with
restriction is multiplication by $[F:\mathbb Q]$, which
is invertible in $L$.  Equivalently the finite-group
kernel in inflation--restriction vanishes by averaging.
This use of a denominator is solely in a
characteristic-zero vanishing argument and has no
integral length consequence.  Apply it to each
coordinate of the restricted class to obtain (G53.2).
\end{proof}

\begin{proposition}[AC-G53-BK: actual cubic tangent cocycles satisfy the source local conditions]
\label{AC-G53-BK}
Let $S$ be the actual cubic ring of AC-G49 and AC-G54-model and let
$F:S\to\overline{\mathbb Q}_\ell[\varepsilon]/(\varepsilon^2)$
be an $O$-algebra map whose first component $f$ is an
\texttt{IsLocPt}.  Its actual tangent cocycle
\[
 z_F(g)=dF(r^{\rm univ}(g).\mathrm{mat})\,
                   f(r^{\rm univ}(g).\mathrm{mat})^{-1}
                                                               \tag{G53.6}
\]
is continuous and belongs to the source's
\texttt{Af.IsBK} local subspace for
$\operatorname{ad}(f(r^{\rm univ}))$.
This assertion applies to all such dual points,
but does not say their global classes vanish.
\end{proposition}
\begin{proof}
The product rule for $F$ and the differential
$\theta_{\rm Lie}(X)=-JX^tJ$ give the adjoint cocycle
identity for (G53.6), including the component with
$\eta=-1$.  This is the seed-independent calculation
\texttt{tan\_coc} in
\path{UnitaryThreeGroups/TanCocIsBK.lean}.
The complete Noetherian $O$-algebra $S$ is an
\texttt{IsCO} ring.  The exact theorem
\path{PadicHodgeTheory/PtFinFinite.lean}, \texttt{Loc.ptFin},
puts both components of $F(S)$ in a finite extension
$L/\mathbb Q_\ell$.  The generic continuity theorem
\texttt{Loc.ptCont}, as used in \texttt{tan\_cont} in
\texttt{TanCocIsBK.lean}, makes both components continuous
on every adically continuous matrix entry.  Matrix
inversion is continuous on the invertible locus, so
$z_F$ is a continuous finite-dimensional cocycle.
These hypotheses depend on the actual complete ring
and its local point, not on the old one-character seed.

For a rational prime $v\ne\ell$, set
$I_v^0=I_v\cap\ker\overline r_\chi$.
The actual universal representation kills $I_v^0$,
so (G53.6) is zero on it.  The residual image is finite,
and $I_v^0$ therefore has finite index in $I_v$.
The adjoint action also kills $I_v^0$.
Thus the restricted cocycle factors through a finite
group.  Its cohomology class is zero in characteristic
zero, by averaging the cocycle over that group.
Explicitly, if $G$ is the finite quotient and
$s=|G|^{-1}\sum_{h\in G}z_F(h)$, the cocycle equation
gives $z_F(g)=s-\rho(g)s$.
This is the exact \texttt{tanBK\_avg} argument in the
source, now with the actual residual kernel.
No integral averaging or claim about its denominator
is made.

At a coefficient prime covered by AC-G49-inert, that
result identifies the entire local
framed functor and selected prime with the old local
one by a constant matrix $C$ over $O$.
The underlying local ring and prime are unchanged.
Conjugating the actual local dual representation by
$C^{-1}$ therefore gives a dual point of the old
selected prime, including its \texttt{IsLocPt}
condition.  Apply the full dual-point theorem
\path{PadicHodgeTheory/RtDual.lean}, \texttt{Loc.rtDual},
and the local/global dictionary
\path{PadicHodgeTheory/BkAbove.lean}, \texttt{Loc.bkAbove}.
They give a $B_{\rm cris}$ cobounding vector for the
conjugated adjoint cocycle.  Conjugating that vector
back by $C$ gives one for $z_F$.
This is valid for the whole dual representation,
including its nilpotent direction; it does not infer
a dual-point condition from a reduced-point condition.
It proves the remaining local field of
\texttt{Af.IsBK}.

For a split coefficient prime with nontrivial local
$\chi$, use the actual prime $P_*$ of AC-G54 instead.
Property 3 of that theorem makes the finite
$L[\varepsilon]/(\varepsilon^2)$ point killing $P_*$
potentially crystalline of the fixed type.  This
uses the full Artin statement, not only its reduced
specialization.  On a finite local extension that
kills the fixed inertial type, its underlying
linear representation is a crystalline extension
of the specialized representation by itself.
As proved in AC-G47-crystalline, the period
dimension criterion splits this extension after
tensoring with $B_{\rm cris}$; equivalently its
adjoint cocycle is a period-ring coboundary.
Finite-index averaging then descends that
coboundary to the rational decomposition group
in characteristic zero.  The local multiplier
is trivial at these split primes, so no extra
semidirect restriction is needed there.
This proves the same \texttt{Af.IsBK} field.
The argument uses no ordered ordinary-inertia
description of all points of the new component.
\end{proof}

\begin{theorem}[AC-G53: the actual cubic basepoint satisfies \texttt{rigid0}]
\label{AC-G53}
Let $x_0:S\to O$ be the actual cubic basepoint.
For every integer $b\leq0$, one has
\[
 \texttt{RigidWt}\bigl(b,r^{\rm univ},\overline{x}_0\bigr).
                                                               \tag{G53.7}
\]
More precisely, a dual point over $\overline{x}_0$ whose
tangent cocycle has weight $b\leq0$ has zero tangent
cocycle and zero derivative on $S$.
Thus the \texttt{rigid0} field of
\path{UnitaryThreeGroups/Rigid0.lean}, \texttt{PtHyp},
is proved for the actual cubic ring.
Neither \texttt{PtHyp.fixed} nor its \texttt{rigid}
field at other points is asserted.
\end{theorem}
\begin{proof}
For the grading torus $\operatorname{diag}(t,1,t^{-1})$,
the nonpositive adjoint pieces of the basepoint are
\[
 \begin{array}{c|c}
 b&\operatorname{ad}(r_\chi)_b\\ \hline
 -1&
 \bigl(\operatorname{Ind}_{H}^{\Gamma}L(\chi^{-j})
                                      \otimes\psi^{-1}\bigr)(-e)\\
 -2&L(\eta\psi^{-2})(-2e)\\
 0&L\oplus L(\eta)\oplus L(\eta).
 \end{array}                                                   \tag{G53.8}
\]
All pieces with $b<-2$ are zero.
For $b=-1$ this is the integral induced-module
calculation (G52.3) and (G52.10), followed by the
projection formula.  For $b=-2$, the matrix unit
$E_{20}$ has $H$-character $\gamma^{-2}$ and the
action of $c$ is multiplication by $-1$; hence its
character is $\eta\gamma^{-2}$.
On diagonal matrices $H$ acts trivially, and
$c$ sends $(a,b,d)$ to $(-d,-b,-a)$.
Its fixed line is $(a,0,-a)$ and its negative
eigenspace has dimension two, proving the last row.
These are actual stable subrepresentations, so their
coordinate projections preserve the local
Bloch--Kato conditions.

Let $F$ be a dual point over $\overline{x}_0$ with
weight-$b$ cocycle.  Its coefficient field is finite
by \texttt{Loc.ptFin}, as in AC-G53-BK.
That proposition puts its cocycle in the Bloch--Kato
subspace of the relevant row of (G53.8).
If the local period cobounding tensor was initially
written with coefficients in $\overline{\mathbb Q}_\ell$,
adjoin its finitely many scalar coefficients to $L$.
The map from the resulting finite-coefficient tensor
product into the algebraic-closure tensor product is
injective.  Its coboundary identity therefore holds
over that finite field as well.  The inertia
averaging vectors away from $\ell$ already lie in
the coefficient field of the cocycle.
Each row is a finite-image Artin representation
followed by a nonpositive Tate twist.  Lemma
AC-G53-Artin, with $n=e,2e$ or $0$, proves that the
cocycle is a coboundary.  This proves (G53.7) with
the exact definition of \texttt{RigidWt} in
\path{UnitaryThreeGroups/GRepresents.lean}.

We verify the stronger zero statement explicitly,
rather than confusing rigidity modulo conjugation
with a zero tangent vector.
The actual ring is normalized.  Differentiating its
normalization makes every $z_F$ a
\texttt{NormCoc}, by the seed-independent
\texttt{tan\_norm} in
\path{UnitaryThreeGroups/TanDet.lean}.
Write $z_F=\partial X$.
Its value at $c$ being zero gives
$X=-JX^tJ$.  Since $2$ is invertible, this forces
\[
 X=\begin{pmatrix}a&u&0\\v&0&-u\\0&-v&-a\end{pmatrix}.
                                                               \tag{G53.9}
\]
Put $t=\chi(g_1)$, $s=\gamma(g_1)$ and
$\Delta=t^j(s-s^{-1})\in O^\times$.
The derivative of the actual \texttt{zmat}
normalization on this coboundary is the commutator
\[
          [X,\operatorname{diag}(\Delta,0,-\Delta)].
                                                               \tag{G53.10}
\]
Its $(0,1)$ and $(1,0)$ entries are $-\Delta u$ and
$\Delta v$.  Both vanish, hence $u=v=0$.
Thus $X=\operatorname{diag}(a,0,-a)$.
This matrix commutes with every diagonal matrix in
(G53.1) and is fixed by $\theta_{\rm Lie}$.
It is therefore fixed by the full adjoint action,
including the nontrivial $\eta$ component.
Consequently $\partial X=0$.
This is also the specialization of the completely
seed-independent theorem
\path{UnitaryThreeGroups/NormCob.lean},
\texttt{normCoc\_cob\_iff}; the displayed calculation
checks all changed seed hypotheses, especially the
actual unit $\Delta$.

Finally (G53.6) and invertibility of the basepoint
matrices show that the derivative of $F$ vanishes on
every actual matrix entry.  It vanishes on $O$ and
therefore on their $O$-algebra.  That algebra is
adically dense in $S$ by AC-G51.
The derivative is continuous by the finite-coefficient
continuity argument of AC-G53-BK, so it is zero on
$S$.  No rigidity statement at a different point
has entered this argument.
\end{proof}

\begin{corollary}[AC-G53-M1: cotangent spanning conditional on the actual \texttt{M2}]
\label{AC-G53-M1}
Let $\mathcal M_j=(A,\mathcal A,x,Q,\ldots)$ be the
actual complete graded model of AC-G51.  If its
degree-zero finiteness condition
$\texttt{M2}(\mathcal A,Q.\mathrm{ideal})$ holds,
then the entire source condition $\texttt{M1}(Q)$ holds.
In particular its equation families in weights
$-1,-2,0$ satisfy the precise saturated cotangent
spanning assertions \texttt{SpansCot}.
This does not prove \texttt{M2}, a numerical cotangent
length bound, or the remaining point conditions
\texttt{M3}--\texttt{M5}.
\end{corollary}
\begin{proof}
The lower bound $1\leq r$ and the specified
\texttt{SmoothPt} chart are already fields of AC-G51.
We prove the three remaining spanning assertions.
Let $D:A\to\overline{\mathbb Q}_\ell$ be a derivation
at $x$, of a weight $b\in\{-1,-2,0\}$, which kills
$Q.\mathrm{ideal}$.
The actual \texttt{Model.der} field of AC-G51 applies
using precisely the assumed \texttt{M2} and the local
basepoint $\overline{x}_0$.
It produces a dual point $F$ of $S$ over
$\overline{x}_0$ whose derivative composed with
$q_A:A\to S$ is $D$.

Every actual matrix entry has a homogeneous preimage
of its matrix-entry weight in $A$.
Since the inverse basepoint matrix is diagonal with
unit entries, multiplication by it in (G53.6) does
not change a matrix-entry weight.
Thus $z_F$ has weight $b$, with no appeal to the old
formula $\gamma^{\mathrm{wt}}$ for the changed seed.
Theorem AC-G53 makes the derivative of $F$ zero, so
$D=0$.

For completeness this vanishing gives the exact
source spanning convention as follows.
Put $I=\ker x$ and
\[
 I^{[2]}=\{a\in A:\ \exists\,0\ne c\in O,\ ca\in I^2\}.
                                                               \tag{G53.11}
\]
This is \texttt{satSq} in
\path{CommutativeAlgebra/ResMapBijective.lean}.
Because $x$ kills nonzero weights, both $I$ and
$I^{[2]}$ are graded.
If the indicated weight-$b$ equation classes did
not span the weight-$b$ part of $I/I^{[2]}$ after
tensoring with $\operatorname{Frac}(O)$, choose a
nonzero linear functional on that quotient,
vanishing on those classes, and extend it by zero
on the other weight parts.  Composing with
$a\mapsto a-x(a)$ produces a nonzero derivation
$D$ of weight $b$.
The derivation rule follows because
\[
 (ab-x(ab))-x(a)(b-x(b))-x(b)(a-x(a))
                         =(a-x(a))(b-x(b))\in I^2.
\]
It kills the equation generators of weight $b$ by
construction, the other generators by weight, and
therefore their entire ideal by the product rule,
since all generators vanish at $x$.
This contradicts the preceding vanishing of $D$.

Spanning after tensoring is exactly
\texttt{SpansCot}: for every integral cotangent
class some nonzero element of $O$ multiplies it
into the $O$-span of the specified equation
classes, obtained by clearing the finitely many
denominators in its expression.
Together with the already constructed chart and
$1\leq r$, this is the literal \texttt{M1}
definition in
\path{CommutativeAlgebra/OrbitVanishingLocal.lean}.
It is the seed-adjusted proof of
\path{UnitaryThreeGroups/GradedModelSpansCotOfRigidWeights.lean},
\texttt{gradedModel\_spansCot\_of\_rigid\_weights}.
Its \texttt{M2} antecedent has been retained, not
removed by rational cohomological vanishing.
\end{proof}

\paragraph{Scope of the reduction and geometric dictionary.}
The actual strict integral extraction is AC-G52.
At every good coefficient prime in the stated
$\ell\ne2,3,269$ range, the present section closes basepoint
\texttt{PtHyp.rigid0} and proves \texttt{M1} conditional
on the actual \texttt{M2}.  It leaves the finite-support
condition, the fixed-point field and rigidity at
nonbase points, especially the projected mixed
rank-four classes, as separate arithmetic inputs.
The finite local systems and their finite splitting
covers use only the previously audited SGA~1,
Expos\'e~V, \S\S4 and~7 dictionary.
The adjoint weight spaces and cotangent calculation
are the explicit representations and first-order
deformation algebras displayed above.  They do not
assert an additional geometric motive or derive
integral torsion sizes from a rational Hodge-weight
vanishing statement.
